\documentclass[11pt]{amsart}

\usepackage[T1]{fontenc}
\usepackage{lmodern}
\usepackage{microtype}
\usepackage{mathtools,amssymb,amsthm}
\usepackage[hidelinks]{hyperref}

\hypersetup{
  pdftitle={Counting collinear triples, four distinct triangles force at most
seven points, and the regular heptagon}
}

\newtheorem{theorem}{Theorem}
\numberwithin{theorem}{section}
\newtheorem{lemma}[theorem]{Lemma}
\newtheorem{proposition}[theorem]{Proposition}
\newtheorem{corollary}[theorem]{Corollary}
\theoremstyle{remark}
\newtheorem{remark}[theorem]{Remark}
\theoremstyle{definition}
\newtheorem*{convention}{Convention}

\DeclareMathOperator{\rank}{rank}
\newcommand{\Tcount}{T}
\newcommand{\Zs}{\mathbb Z}

\title[Four distinct triangles, collinear triples counted]
{Counting Collinear Triples, Four Distinct Triangles Force at Most Seven
Points, and the Regular Heptagon}

\date{}

\subjclass[2020]{52C10, 51M04, 52C35}
\keywords{distinct triangles, few-distance sets, regular heptagon,
Cayley--Menger determinant, exhaustive census, Schoenberg's criterion}

\begin{document}

\begin{abstract}
Count a triangle as a triple of distinct points of the plane up to rigid motion,
collinear triples included, as Palsson and Yu do. With triangles counted that
way, and not under the rival convention noted below, we prove their
Conjecture~4.1: a planar set spanning at most four distinct triangles has at
most seven points, and a seven-point set spanning four distinct triangles is
similar to the vertex set of a regular heptagon. The proof
is a complete census of the $161\,351$ edge-colourings of $K_7$ that use at
least three distance values and span at most four triangles, an elimination of
all but one $S_7$-orbit of them by exact polynomial certificates in integer and
rational arithmetic, a Fourier computation identifying the surviving circulant
as the regular heptagon, and a short hand argument for eight points and more.
Under the rival convention that counts only non-collinear triples neither half
of the conjecture is decided here, and the uniqueness half is false: the regular
hexagon with its centre is a seven-point counterexample.
\end{abstract}

\maketitle

\section{Introduction}

For a finite set $S\subset\mathbb R^2$ let $\Tcount(S)$ be the number of
equivalence classes of $3$-element subsets of $S$ under rigid motions of the
plane. Palsson and Yu \cite{PY} fix this counting convention explicitly; we
quote it verbatim:

\begin{convention}
``To avoid ambiguity regarding degenerate triangles, we can formalize this as
counting the number of triples of distinct points in a point set, up to the
equivalence relation defined by rigid plane isometries (hence, the three
vertices of the triangle are allowed to be collinear, but must be distinct).''
\end{convention}

\noindent
So collinear triples are counted, and by side--side--side a class is exactly a
multiset of three pairwise distances. Their Theorem~1.1 states that a planar set
spanning at most three distinct triangles has at most six points, uniquely the
regular hexagon, and their Section~4 then reads, verbatim:

\begin{quote}
``Firstly, observing that the regular heptagon has $4$ distinct triangles, we
conjecture that this is optimal:''\\[2pt]
\textbf{Conjecture 4.1.} ``The maximum number of points in the plane spanning
four distinct triangles is $7$, achieved only by the regular heptagon.''
\end{quote}

\begin{theorem}\label{thm:main}
Count triangles as in the Convention. Then
\begin{enumerate}
\item[(a)] every set of seven points in the plane spans at least four distinct
  triangles;
\item[(b)] a set of seven points in the plane spanning at most four distinct
  triangles is similar to the vertex set of a regular heptagon;
\item[(c)] no set of eight or more points in the plane spans at most four
  distinct triangles.
\end{enumerate}
In particular $\max\{|S| : \Tcount(S)=4\}=7$, attained only by the regular
heptagon: Conjecture~4.1 of \cite{PY} holds with triangles counted as in the
Convention. It is not decided here under the rival convention described in the
Scope paragraph below.
\end{theorem}

Part~(a) reproves the seven-point case of \cite[Theorem~1.1]{PY} and is used in
the proof of~(c). Parts~(a) and~(b) are Proposition~\ref{prop:seven} below;
part~(c) is Proposition~\ref{prop:eight}.

\medskip
\noindent\textbf{Scope.}
Theorem~\ref{thm:main} is proved under the Convention, the count \cite{PY}
states. Under the rival convention that counts only non-collinear triples --- used by
\cite[Definitions 3.2--3.3]{ELMP} and by \cite{BDPS} --- this paper decides
\emph{neither} half of Conjecture~4.1, and the uniqueness half is then
\emph{false}: the regular hexagon together with its centre is a seven-point set
with exactly four non-collinear classes which is not similar to a regular
heptagon (Remark~\ref{rem:hexcentre}). Our census enumerates the cell ``at most
four classes \emph{with} collinear triples counted'', and that configuration has
five classes in our count, so it was never in the cell; for the same reason
nothing here bears on the maximality half under the rival convention, which we
leave open, nor on \cite[Conjecture~2.1]{ELMP} (``any set of seven points in the
plane determines at least four distinct triangles''), which is stated under the
non-collinear count.

Lemmas \ref{lem:census}--\ref{lem:small} are finite computations, carried out in
exact integer and rational arithmetic by the program accompanying this note;
Sections \ref{sec:objects}, \ref{sec:certs} and~\ref{sec:eight} are by hand, and
Section~\ref{sec:verify} says what the program covers. The census splits by the
number $D$ of distinct distances, and $156\,710$ of its $161\,351$ leaves have
$D=3$; that lane is covered by Shinohara's classification of planar
three-distance sets \cite{Shinohara}, which is not used here. Enumerate
colourings, then eliminate by realizability, is the method of Sz\H{o}ll\H{o}si
and \"Osterg\aa rd \cite{SO}.

\section{The two configurations}\label{sec:objects}

Let $\zeta=e^{2\pi i/7}$ and let $H=\{\zeta^0,\dots,\zeta^6\}$ be the regular
heptagon. For $g\in\{1,2,3\}$ put
\[
  s_g \;=\; |1-\zeta^g|^2 \;=\; 2-2\cos(2\pi g/7),
\]
so $s_1,s_2,s_3$ are the three squared chord lengths. With $c=2\cos(2\pi/7)$,
which satisfies $c^3+c^2-2c-1=0$, one has $s_1=2-c$, $s_2=4-c^2$ and
$s_3=1+c+c^2$, and
\begin{equation}\label{eq:chords}
  s_1,s_2,s_3 \text{ are the three roots of } x^3-7x^2+14x-7 .
\end{equation}
All three are positive, because for $x\le 0$ every term of $x^3-7x^2+14x-7$ is
$\le 0$ and its constant term is $-7<0$; and they are distinct, that cubic being
irreducible over $\mathbb Q$. Numerically
$s_1=0.75302\ldots$, $s_2=2.44504\ldots$, $s_3=3.80193\ldots$.

\begin{lemma}\label{lem:heptagon}
$\Tcount(H)=4$ and $H$ realises exactly three distinct distances. No triple of
vertices of $H$ is collinear, so $\Tcount(H)=4$ under either convention.
\end{lemma}

\begin{proof}
The chord $\zeta^i\zeta^j$ has squared length $s_g$ with
$g=\min(|i-j|,7-|i-j|)$, and $g\mapsto s_g$ is injective by \eqref{eq:chords},
so a triple's class is the multiset of the three cyclic gaps it cuts out. Those
gaps are a partition of $7$ into three positive parts read modulo the folding
$g\mapsto\min(g,7-g)$, and the four partitions $5{+}1{+}1$, $4{+}2{+}1$,
$3{+}3{+}1$, $3{+}2{+}2$ give the four multisets
\[
  \{s_1,s_1,s_2\},\quad \{s_1,s_2,s_3\},\quad \{s_1,s_3,s_3\},\quad
  \{s_2,s_2,s_3\},
\]
and no others; all $\binom73=35$ triples are accounted for. A triple with
squared side lengths $a,b,c$ is collinear exactly when
$2(ab+bc+ca)-(a^2+b^2+c^2)=0$, that expression being sixteen times the squared
area; it is a nonzero element of $\mathbb Q(c)$ for each of the four multisets
above.
\end{proof}

\begin{remark}\label{rem:hexcentre}
Let $X$ consist of the six vertices of a unit regular hexagon together with its
centre, say
\[
  (\pm 1,0),\ \left(\pm\tfrac12,\pm\tfrac{\sqrt3}2\right),\ (0,0).
\]
Its $21$ squared distances take the three integer values $1,3,4$, and its $35$
triples fall into the five classes
$\{1,1,1\}$, $\{1,1,3\}$, $\{1,1,4\}$, $\{1,3,4\}$, $\{3,3,3\}$. Exactly one of
them, $\{1,1,4\}$, is collinear --- a vertex, the centre and the opposite vertex
--- and it occurs for three triples, whereas $H$ has none. So $\Tcount(X)=5$
under the Convention, while $X$ spans only four \emph{non-collinear} classes on
seven points and is not similar to $H$.
\end{remark}

\section{Colourings, realizability, and two certificates}\label{sec:certs}

Let $S=\{p_0,\dots,p_{n-1}\}\subset\mathbb R^2$ consist of $n$ distinct points
and let $d_{ij}=|p_i-p_j|^2$. Colour the edge $\{i,j\}$ of $K_n$ by the
\emph{value} $d_{ij}$; write $D$ for the number of colours. By
side--side--side, two triples are congruent if and only if their multisets of
edge colours agree, so
\begin{equation}\label{eq:Tequals}
  \Tcount(S) \;=\; \#\{\text{distinct multisets of edge colours of the
  $\binom n3$ triangles of } K_n\}.
\end{equation}
Number the edges $(0,1),(0,2),(1,2),(0,3),\dots,(n-2,n-1)$ and number the
colours $1,2,\dots$ in order of first appearance; this puts the colouring in a
unique \emph{restricted-growth word} $w(S)\in\{1,\dots,D\}^{\binom n2}$, and it
quotients the relabelling of colours exactly. Scaling $S$ scales all $d_{ij}$
together, so we may and do normalise the squared distance of colour~$1$ to $1$
and write $y_2,\dots,y_D>0$ for the others; these are pairwise distinct and
distinct from $1$, since distinct colours are distinct distance values.

We use one classical criterion and two elementary consequences of it.

\begin{lemma}[Schoenberg \cite{Schoenberg}]\label{lem:schoenberg}
Let $M=(m_{ij})$ be a real symmetric $n\times n$ matrix with zero diagonal and
fix $b$. There are points $p_0,\dots,p_{n-1}\in\mathbb R^k$ with
$|p_i-p_j|^2=m_{ij}$ for all $i,j$ if and only if the $(n-1)\times(n-1)$ matrix
\[
  G^{(b)} = \left(\tfrac12\left(m_{bi}+m_{bj}-m_{ij}\right)\right)_{i,j\ne b}
\]
is positive semidefinite of rank at most $k$. \textup{(}$G^{(b)}$ is the Gram
matrix of the vectors $p_i-p_b$.\textup{)}
\end{lemma}

We shall use the criterion in the following equivalent form. Let
$H_0=\{x\in\mathbb R^n:\sum_ix_i=0\}$. If $m_{ij}=|p_i-p_j|^2$ and
$\sum_ix_i=0$ then
\[
  x^\top Mx=\sum_{i,j}x_ix_j\bigl(|p_i|^2-2\langle p_i,p_j\rangle+|p_j|^2\bigr)
  =-2\Bigl|\sum_i x_ip_i\Bigr|^2 ,
\]
and $\sum_ix_ip_i=\sum_{i\ne b}x_i(p_i-p_b)$, so under
$x\mapsto(x_i)_{i\ne b}$ the form $-M$ restricted to $H_0$ is twice the
quadratic form of $G^{(b)}$. In particular $-M|_{H_0}$ is positive semidefinite
of rank exactly $\rank G^{(b)}$, so realizability in $\mathbb R^k$ says: $-M$ is
positive semidefinite on $H_0$ with at most $k$ nonzero eigenvalues there.

\begin{corollary}\label{cor:minors}
If $S\subset\mathbb R^2$ then, for every $b$, every $3\times3$ minor of
$G^{(b)}$ vanishes; equivalently, every $4$-point Cayley--Menger determinant
\[
  \mathrm{CM}(i,j,k,l)=
  \det\begin{pmatrix}
    0&1&1&1&1\\ 1&0&d_{ij}&d_{ik}&d_{il}\\
    1&d_{ij}&0&d_{jk}&d_{jl}\\ 1&d_{ik}&d_{jk}&0&d_{kl}\\
    1&d_{il}&d_{jl}&d_{kl}&0
  \end{pmatrix}
\]
vanishes.
\end{corollary}

Fix a restricted-growth word $w$ on $K_n$. Each $\mathrm{CM}(i,j,k,l)$ and each
$3\times3$ minor of each $G^{(b)}$ is then a polynomial with integer
coefficients in $y_2,\dots,y_D$, computable from $w$ alone; call these the
\emph{test polynomials} of $w$. By Corollary~\ref{cor:minors} all of them
vanish at any planar realization of $w$, whose coordinates $y$ lie in the open
positive orthant. Hence:

\begin{lemma}[one-signed certificate]\label{lem:onesigned}
If some test polynomial of $w$ is nonzero and all of its coefficients have the
same sign, then $w$ has no planar realization.
\end{lemma}

\begin{proof}
All exponents are non-negative, so such a polynomial is strictly positive, or
strictly negative, at every point of the positive orthant, and in particular
does not vanish there.
\end{proof}

\begin{lemma}[univariate gcd certificate]\label{lem:gcd}
Let $U$ be the set of test polynomials of $w$ that involve the single variable
$y_i$, and let $g=\gcd U\in\mathbb Q[y_i]$, taken monic. Then $g(y_i)=0$ at any
planar realization. Consequently, writing $g=y_i^{\,k}h$ with $h(0)\ne0$:
\begin{enumerate}
\item[(i)] if $h$ is constant and $k=0$, then $w$ has no planar realization;
\item[(ii)] if $h$ is constant and $k\ge1$, then $y_i=0$, which is impossible
  for distinct points, so $w$ has no planar realization;
\item[(iii)] if $\deg h=1$ and its unique root $r$ is positive, then $y_i=r$,
  and substituting $y_i=r$ into every test polynomial of $w$ leaves a system in
  the remaining variables to which (i)--(iii) may be applied again.
\end{enumerate}
\end{lemma}

\begin{proof}
Every element of $U$ vanishes at a realization, hence so does any polynomial
combination of them, hence so does $g$. In~(iii), $y_i>0$ forces $y_i=r$ since
$r$ is the only root of $h$ and $y_i^k\ne0$.
\end{proof}

\section{Seven points}\label{sec:seven}

Throughout this section $n=7$, so $K_7$ has $21$ edges and
$\binom73=35$ triples, and the \emph{cell} is the set of restricted-growth words
on $K_7$ with at most four distinct triple colour-multisets: by
\eqref{eq:Tequals} the word of any seven-point set with $\Tcount(S)\le4$ lies in
the cell.

\begin{lemma}\label{lem:census}
Exactly $161\,351$ words of the cell use $D\ge3$ colours. Every one of them has
$D\le4$; $156\,710$ have $D=3$ and $4\,641$ have $D=4$. Split further by the
number of triple multisets realised, the counts are $3\,948$ with $(D,\Tcount)=(3,3)$,
$152\,762$ with $(3,4)$ and $4\,641$ with $(4,4)$.
\end{lemma}

\begin{proof}
A depth-first traversal that assigns colours to the $21$ edges in order,
maintaining the set of triple multisets already completed and pruning as soon as
a fifth one appears. The traversal is complete and uncapped: it explores
$2\,915\,582$ nodes and terminates. That $3\,948+152\,762+4\,641=161\,351$ is
the arithmetic of the exhaustion.
\end{proof}

\begin{lemma}\label{lem:filter}
Of those $161\,351$ words, $81\,306$ have a $4$-subset whose Cayley--Menger
determinant is a nonzero constant, and $70\,286$ further ones have a $4$-subset
whose Cayley--Menger determinant is non-constant with all coefficients of one
sign. By Lemma~\ref{lem:onesigned} none of these $151\,592$ words is
realizable. The remaining $9\,759$ words form exactly $30$ orbits under the
action of $S_7$ on the vertices.
\end{lemma}

\begin{proof}
A finite computation, $35$ determinants per word; $81\,306+70\,286+9\,759=161\,351$.
The test depends on the word only through its $4$-subsets, so the surviving set
is $S_7$-stable, and the orbit count is a union--find over two generators of
$S_7$. The orbit sizes, in the notation
$\text{size}^{\,\text{multiplicity}}$, are
$70$, $105^{\,5}$, $120$, $140$, $210^{\,5}$, $252^{\,2}$, $315^{\,8}$,
$420$, $630^{\,5}$, $1260$, and they sum to $9\,759$.
\end{proof}

\begin{lemma}\label{lem:elim}
Of those $30$ orbits, $27$ carry a $3\times3$ Gram minor that is nonzero with
all coefficients of one sign, and two more are killed by
Lemma~\ref{lem:gcd}. Exactly one orbit remains, the orbit of size $120$ of the
word
\begin{equation}\label{eq:W}
  W=1\,1\,2\,2\,1\,3\,2\,3\,1\,3\,3\,2\,3\,1\,2\,3\,3\,2\,2\,1\,1 ,
\end{equation}
listed against the edge order
\begin{gather*}
  01,\;02,\;12,\;03,\;13,\;23,\;04,\;14,\;24,\;34,\;05,\\
  15,\;25,\;35,\;45,\;06,\;16,\;26,\;36,\;46,\;56 .
\end{gather*}
\end{lemma}

\begin{proof}
For $27$ orbits, Lemma~\ref{lem:onesigned}. The two others have $D=4$, and for
each of them we record the monic gcd's of Lemma~\ref{lem:gcd} in the order the
substitutions are made. For
$1\,1\,2\,1\,3\,4\,1\,4\,3\,2\,2\,1\,1\,1\,1\,3\,1\,1\,1\,1\,4$
(orbit size $210$): the $8$ test polynomials in $y_2$ alone have gcd
$y_2^3-2y_2^2$, so $y_2=2$; then the $32$ in $y_3$ alone have gcd
$y_3^2-2y_3$, so $y_3=2$; then the $2\,192$ in $y_4$ alone have gcd $y_4$,
forcing $y_4=0$, which case~(ii) of Lemma~\ref{lem:gcd} excludes. For
$1\,1\,2\,2\,1\,1\,3\,4\,4\,3\,3\,4\,4\,3\,2\,4\,3\,3\,4\,1\,1$
(orbit size $105$): gcd $y_2^3-2y_2^2$ gives $y_2=2$; gcd $y_3-1$ gives
$y_3=1$; and then the $2\,288$ test polynomials in $y_4$ alone have gcd $1$, so
by case~(i) they have no common root at all. (The value $y_3=1$ already
contradicts $y_3\ne1$, so that orbit dies twice over.) Finally $27+2+1=30$.
\end{proof}

\begin{lemma}\label{lem:heptagononly}
A seven-point planar set whose word is $W$ is similar to $H$.
\end{lemma}

\begin{proof}
Relabelling the vertices by $\sigma=(0,1,6,2,5,3,4)$, meaning $0\mapsto0$,
$1\mapsto1$, $2\mapsto6$, $3\mapsto2$, $4\mapsto5$, $5\mapsto3$,
$6\mapsto4$, carries $W$ to the \emph{gap colouring} of $\Zs_7$, in which the
colour of $\{i,j\}$ is $\min(|i-j|,7-|i-j|)$; there are $42=|\mathrm{AGL}(1,7)|$
such relabellings, consistent with the orbit size $5040/42=120$. Each colour
class of $W$ is a single $7$-cycle; colour~1, for instance, is
$0\,1\,3\,5\,6\,4\,2$.

So the squared-distance matrix is the circulant $M$ on $\Zs_7$ with
$M_{ij}=y_{g}$, $g=\min(|i-j|,7-|i-j|)$, for positive $y_1,y_2,y_3$. Its
eigenvectors are the Fourier modes $v_m=(\omega^{jm})_j$, $\omega=e^{2\pi i/7}$,
with
\[
  \mu_m \;=\; \sum_{g=1}^{3} y_g\,c_{gm},
  \qquad c_k := 2\cos(2\pi k/7)=c_{7-k},
\]
so that $v_0=(1,\dots,1)$ has eigenvalue $\mu_0=2(y_1+y_2+y_3)>0$. The modes
$v_1,\dots,v_6$ span $H_0$, and $\mu_m=\mu_{7-m}$, so
Lemma~\ref{lem:schoenberg}, in the equivalent form derived just after it, says
that realizability in $\mathbb R^2$ amounts to $\mu_m\le0$
for $m=1,\dots,6$ with at most two of them nonzero. Nonzero eigenvalues come
in equal pairs, so their number is even, hence $0$ or $2$; it is not $0$, since
$M\ne0$. Therefore exactly one of $\mu_1,\mu_2,\mu_3$ is negative and the other
two vanish.

Say $\mu_2=\mu_3=0$. These are two linear equations in $(y_1,y_2,y_3)$ over
$\mathbb Q(c)$, and the $2\times2$ minor $c_2c_1-c_3c_3$ of their coefficient
matrix is a nonzero element of that field, so their solution space is a single
line. The vector $(s_1,s_2,s_3)$ lies on it and gives $\mu_1=-7$, both by direct
computation in $\mathbb Q(c)$; hence it \emph{spans} it, and
$(y_1,y_2,y_3)=\lambda(s_1,s_2,s_3)$ with $\lambda>0$. By
Lemma~\ref{lem:schoenberg} the point set is then determined up to isometry, and
it is $\sqrt\lambda\,H$.

The two remaining cases are the same set relabelled. The multiplier
$j\mapsto2j$ of $\Zs_7$ sends gap $g$ to gap $2g$, hence permutes $(1,2,3)$ as
$(2,3,1)$, and $j\mapsto3j$ as $(3,1,2)$; correspondingly
$(y_1,y_2,y_3)\propto(s_2,s_3,s_1)$ gives $(\mu_1,\mu_2,\mu_3)=(0,-7,0)$ and
$(s_3,s_1,s_2)$ gives $(0,0,-7)$, again by computation in $\mathbb Q(c)$. Both
are the regular heptagon with its three distance classes renamed.
\end{proof}

\begin{lemma}\label{lem:small}
No seven-point planar set has $D\le2$.
\end{lemma}

\begin{proof}
$D=1$ is impossible already for four points: the Cayley--Menger determinant of a
monochromatic $4$-subset is a nonzero constant. For $D=2$ it suffices, by
heredity, to rule out six points. There are $2^{14}-1=16\,383$
restricted-growth words on the $15$ edges of $K_6$ using exactly two colours,
and all of them lie in the cell, two colours yielding at most four triple
multisets automatically. The Cayley--Menger test of Lemma~\ref{lem:onesigned}
leaves $11\,242$ of them, forming $42$ orbits under $S_6$; of these, $33$ carry
a one-signed $3\times3$ Gram minor and the remaining $9$ are killed by
Lemma~\ref{lem:gcd}. So no planar six-point two-distance set exists, which is
the upper half of Erd\H{o}s and Fishburn's $g(2)=5$ \cite{EF}.
\end{proof}

\begin{proposition}\label{prop:seven}
Every seven-point planar set spans at least four distinct triangles, and one
spanning at most four is similar to $H$.
\end{proposition}

\begin{proof}
Let $|S|=7$ with $\Tcount(S)\le4$. By Lemma~\ref{lem:small}, $S$ has $D\ge3$
distinct distances, so by \eqref{eq:Tequals} its word is one of the $161\,351$
words of Lemma~\ref{lem:census}; by Lemmas \ref{lem:filter} and~\ref{lem:elim}
that word lies in the orbit of $W$; by Lemma~\ref{lem:heptagononly} $S$ is
similar to $H$. That gives~(b), and $\Tcount(H)=4$ by
Lemma~\ref{lem:heptagon} gives~(a): a seven-point set with
$\Tcount(S)\le3$ would be similar to $H$ and so have $\Tcount(S)=4$, a
contradiction.
\end{proof}

\section{Eight points and more}\label{sec:eight}

\begin{proposition}\label{prop:eight}
No planar set of eight or more points spans at most four distinct triangles.
\end{proposition}

\begin{proof}
Suppose $|S|=8$ and $\Tcount(S)\le4$. For $p\in S$ the classes of $S\setminus\{p\}$
are among those of $S$, so $\Tcount(S\setminus\{p\})\le4$; and
$\Tcount(S\setminus\{p\})\ge4$ by Proposition~\ref{prop:seven}(a). Hence
$\Tcount(S\setminus\{p\})=4$ and, by Proposition~\ref{prop:seven}(b),
$S\setminus\{p\}$ is the vertex set of a regular heptagon, for every
$p\in S$.

Pick $p\ne q$ in $S$ and put $R=S\setminus\{p,q\}$, so $|R|=6$ and
$R\cup\{q\}$, $R\cup\{p\}$ are both regular heptagons. Six vertices of a
regular heptagon determine the seventh. Indeed they are concyclic and no three
are collinear, so the circle through any three of them is the circumcircle and
is determined by $R$; the six points cut that circle into six arcs, of which
exactly one has length $2\cdot(2\pi/7)$ and the others $2\pi/7$; and the
missing vertex is the midpoint of the long arc. Applying this to $R$ inside
$R\cup\{p\}$ and inside $R\cup\{q\}$ gives $p=q$, a contradiction.

If $|S|\ge9$ and $\Tcount(S)\le4$, any eight-point subset $S'$ has
$\Tcount(S')\le4$ by the same heredity, contradicting the previous paragraph.
\end{proof}

Theorem~\ref{thm:main} follows: (a) and (b) are Proposition~\ref{prop:seven},
(c) is Proposition~\ref{prop:eight}, and $\Tcount(H)=4$ with $|H|=7$ shows the
maximum is attained.

\section{Verification}\label{sec:verify}

The objects the hand arguments use are printed above: the heptagon and its four
classes in Lemma~\ref{lem:heptagon}, the hexagon-plus-centre in
Remark~\ref{rem:hexcentre}, the surviving word $W$ in \eqref{eq:W} with its edge
order and the relabelling $\sigma$ that turns it into the gap colouring of
$\Zs_7$, the two eliminated $D=4$ words with their gcd chains in
Lemma~\ref{lem:elim}, and the squared chords in \eqref{eq:chords}. The census
itself, the $161\,351$ words and the Cayley--Menger counts of
Lemma~\ref{lem:filter}, the $27$ one-signed Gram minors of
Lemma~\ref{lem:elim} and both lanes of Lemma~\ref{lem:small} are not printed
here; they are supplied through the program accompanying this note.

That program recomputes, in exact integer and \texttt{Fraction} arithmetic and
using the Python standard library only, every numerical claim above: the squared
chords and their minimal polynomial and the four heptagon classes with their
non-degeneracy, over $\mathbb Q[c]/(c^3+c^2-2c-1)$; the five classes and one
collinear class of the hexagon-plus-centre, over $\mathbb Q$; the census of
Lemma~\ref{lem:census} including its node count and its cross-tabulation; the
three counts, the $30$ orbits and their size profile in
Lemma~\ref{lem:filter}; the $27+2+1$ split of Lemma~\ref{lem:elim} with the gcd
chains printed; the relabelling $\sigma$, the eigenvalue vectors $(-7,0,0)$,
$(0,-7,0)$, $(0,0,-7)$ and the rank-two minor of
Lemma~\ref{lem:heptagononly}; and both lanes of Lemma~\ref{lem:small}. It prints
one line per check and exits $0$ only if all of them pass; all $48$ passed.

Nothing else is machine-checked. Schoenberg's criterion (Lemma~\ref{lem:schoenberg}) is
quoted, not proved; Proposition~\ref{prop:eight} contains no computation; and
beyond the single configuration of Remark~\ref{rem:hexcentre} nothing about the
non-collinear convention is checked. The two certificates are sufficient
conditions for non-realizability, and no claim is made that the $30$ orbits
could not be eliminated another way.

\begin{thebibliography}{9}

\bibitem{BDPS}
H.~N. Brenner, J.~S. Depret-Guillaume, E.~A. Palsson, and S.~Senger,
\emph{Uniqueness of optimal point sets determining two distinct triangles},
Integers \textbf{21} (2021), Paper No.~A43;
\href{https://arxiv.org/abs/1910.00629}{arXiv:1910.00629}.
Its triangle count is the non-collinear one.

\bibitem{ELMP}
A.~Epstein, A.~Lott, S.~J. Miller, and E.~A. Palsson,
\emph{Optimal point sets determining few distinct triangles},
Integers \textbf{18} (2018), Paper No.~A16.
Definitions 3.2--3.3 use the non-collinear count; Conjecture~2.1 is stated
there.

\bibitem{EF}
P.~Erd\H{o}s and P.~Fishburn,
\emph{Maximum planar sets that determine $k$ distances},
Discrete Math. \textbf{160} (1996), 115--125.

\bibitem{PY}
E.~A. Palsson and E.~Yu,
\emph{On optimal point sets determining distinct triangles},
Electron. J. Combin. \textbf{31} (2024), no.~2, Paper No.~P2.24.
\href{https://doi.org/10.37236/12360}{doi:10.37236/12360};
\href{https://arxiv.org/abs/2308.13107}{arXiv:2308.13107}.
Quotations above are verbatim from the source of arXiv:2308.13107v2.

\bibitem{Schoenberg}
I.~J. Schoenberg,
\emph{Remarks to Maurice Fr\'echet's article ``Sur la d\'efinition axiomatique
d'une classe d'espaces distanci\'es vectoriellement applicable sur l'espace de
Hilbert''},
Ann. of Math. (2) \textbf{36} (1935), 724--732.

\bibitem{Shinohara}
M.~Shinohara,
\emph{Classification of three-distance sets in two dimensional Euclidean
space},
European J. Combin. \textbf{25} (2004), 1039--1058.
\href{https://doi.org/10.1016/j.ejc.2003.12.009}{doi:10.1016/j.ejc.2003.12.009}.

\bibitem{SO}
F.~Sz\H{o}ll\H{o}si and P.~R.~J. \"Osterg\aa rd,
\emph{Constructions of maximum few-distance sets in Euclidean spaces},
Electron. J. Combin. \textbf{27} (2020), no.~1, Paper No.~P1.23.
\href{https://doi.org/10.37236/8565}{doi:10.37236/8565}.

\end{thebibliography}

\end{document}
